% GENERATED FILE. Edit canonical lessons, then run npm run generate:books.
% canonical-source-hash: fnv1a64:1ea3141bf34e354b
% canonical-lessons: JA-C79-mukou, JA-C79-temae, JA-C79-kado, JA-C79-kaidan, JA-C79-rouka

\chapter{Other side, This side of, Corner, Stairs, Corridor}
\label{ch:ja-other-side-this-side-of-corner-stairs-corridor}
\hlchaptermodality{79}

\begin{chapteropening}
\textbf{By the end of this chapter:} \emph{I can say the other side, this side of, a corner, stairs and a corridor.}

Complete the last lesson of chapter 79: I can say the other side, this side of, a corner, stairs and a corridor.
\end{chapteropening}

\section[\texorpdfstring{\ja{むこう} (mukou)}{むこう (mukou)}]{\ja{むこう} (mukou) — the other side}

\label{lesson:JA-C79-mukou}



\begin{quote}
Before the new one: say the Japanese for beside, then the Japanese for the very middle.
\end{quote}

\subsection*{You'll want to know: \ja{むこう}}
\textbf{\ja{むこう}} — \emph{mukou} — \textquotedblleft{}the other side\textquotedblright{}.

\begin{grammarlens}[title={What you've built}]
One more word for this chapter.
\end{grammarlens}

\subsection*{Your turn}
\begin{itemize}
\raggedright
  \item \emph{Say it:} \emph{mukou}
  \item \emph{Say it:} \emph{mukou}, once more
  \item \emph{Say it:} say \emph{mukou}
  \item \emph{Recall:} say the Japanese for beside, then the Japanese for the very middle, then say \emph{mukou} again
\end{itemize}

\begin{tcolorbox}[breakable,colback=teal!4,colframe=teal!35!black,title={Before you move on}]
What does \ja{むこう} mean? (\textquotedblleft{}The other side\textquotedblright{}.) Say it once more.
\end{tcolorbox}
\section[\texorpdfstring{\ja{てまえ} (temae)}{てまえ (temae)}]{\ja{てまえ} (temae) — this side of}

\label{lesson:JA-C79-temae}



\begin{quote}
Before the new one: say the Japanese for the very middle, then the Japanese for the other side.
\end{quote}

\subsection*{You'll want to know: \ja{てまえ}}
\textbf{\ja{てまえ}} — \emph{temae} — \textquotedblleft{}this side of\textquotedblright{}.

\begin{grammarlens}[title={What you've built}]
One more word for this chapter.
\end{grammarlens}

\subsection*{Your turn}
\begin{itemize}
\raggedright
  \item \emph{Say it:} \emph{temae}
  \item \emph{Say it:} \emph{temae}, once more
  \item \emph{Say it:} say \emph{temae}
  \item \emph{Recall:} say the Japanese for the very middle, then the Japanese for the other side, then say \emph{temae} again
\end{itemize}

\begin{tcolorbox}[breakable,colback=teal!4,colframe=teal!35!black,title={Before you move on}]
What does \ja{てまえ} mean? (\textquotedblleft{}This side of\textquotedblright{}.) Say it once more.
\end{tcolorbox}
\section[\texorpdfstring{\ja{かど} (kado)}{かど (kado)}]{\ja{かど} (kado) — a corner}

\label{lesson:JA-C79-kado}



\begin{quote}
Before the new one: say the Japanese for the other side, then the Japanese for this side of.
\end{quote}

\subsection*{You'll want to know: \ja{かど}}
\textbf{\ja{かど}} — \emph{kado} — \textquotedblleft{}a corner\textquotedblright{}.

\begin{grammarlens}[title={What you've built}]
One more word for this chapter.
\end{grammarlens}

\subsection*{Your turn}
\begin{itemize}
\raggedright
  \item \emph{Say it:} \emph{kado}
  \item \emph{Say it:} \emph{kado}, once more
  \item \emph{Say it:} say \emph{kado}
  \item \emph{Recall:} say the Japanese for the other side, then the Japanese for this side of, then say \emph{kado} again
\end{itemize}

\begin{tcolorbox}[breakable,colback=teal!4,colframe=teal!35!black,title={Before you move on}]
What does \ja{かど} mean? (\textquotedblleft{}A corner\textquotedblright{}.) Say it once more.
\end{tcolorbox}
\section[\texorpdfstring{\ja{かいだん} (kaidan)}{かいだん (kaidan)}]{\ja{かいだん} (kaidan) — stairs}

\label{lesson:JA-C79-kaidan}



\begin{quote}
Before the new one: say the Japanese for this side of, then the Japanese for a corner.
\end{quote}

\subsection*{You'll want to know: \ja{かいだん}}
\textbf{\ja{かいだん}} — \emph{kaidan} — \textquotedblleft{}stairs\textquotedblright{}.

\begin{grammarlens}[title={What you've built}]
One more word for this chapter.
\end{grammarlens}

\subsection*{Your turn}
\begin{itemize}
\raggedright
  \item \emph{Say it:} \emph{kaidan}
  \item \emph{Say it:} \emph{kaidan}, once more
  \item \emph{Say it:} say \emph{kaidan}
  \item \emph{Recall:} say the Japanese for this side of, then the Japanese for a corner, then say \emph{kaidan} again
\end{itemize}

\begin{tcolorbox}[breakable,colback=teal!4,colframe=teal!35!black,title={Before you move on}]
What does \ja{かいだん} mean? (\textquotedblleft{}Stairs\textquotedblright{}.) Say it once more.
\end{tcolorbox}
\section[\texorpdfstring{\ja{ろうか} (rouka)}{ろうか (rouka)}]{\ja{ろうか} (rouka) — a corridor}

\label{lesson:JA-C79-rouka}



\begin{quote}
Before the new one: say the Japanese for a corner, then the Japanese for stairs.
\end{quote}

\subsection*{You'll want to know: \ja{ろうか}}
\textbf{\ja{ろうか}} — \emph{rouka} — \textquotedblleft{}a corridor\textquotedblright{}.

\begin{grammarlens}[title={What you've built}]
One more word for this chapter.
\end{grammarlens}

\subsection*{Your turn}
\begin{itemize}
\raggedright
  \item \emph{Say it:} \emph{rouka}
  \item \emph{Say it:} \emph{rouka}, once more
  \item \emph{Say it:} say \emph{rouka}
  \item \emph{Recall:} say the Japanese for a corner, then the Japanese for stairs, then say \emph{rouka} again
\end{itemize}

\begin{tcolorbox}[breakable,colback=teal!4,colframe=teal!35!black,title={Before you move on}]
What does \ja{ろうか} mean? (\textquotedblleft{}A corridor\textquotedblright{}.) Say it once more.
\end{tcolorbox}
